\setlength{\footskip}{8mm}

\chapter{Introduction}
\label{ch:introduction}

\section{Overview}
\label{sec:intro-overview}

Human pose estimation from monocular video supports a wide range of
applications, including motion capture for animation, human-computer
interaction, sports performance analysis, surveillance, and autonomous
driving \cite{guo2025survey3dhpe}. Among the various formulations of this
problem, monocular 3D human pose estimation (HPE), which recovers the
three-dimensional coordinates of human body joints from a single-camera
video sequence, remains particularly challenging, since a single 2D
projection is inherently ambiguous with respect to depth. The dominant
paradigm decomposes the task into two stages: a 2D pose detector localizes
keypoints on the image plane, after which a 2D-to-3D lifting network infers
the corresponding 3D joint positions from the resulting 2D keypoint
sequence \cite{pavllo2019videopose3d}. Depth ambiguity and self-occlusion
make this second stage, lifting, the central technical difficulty the
field has focused on for the past decade.

Early lifting approaches relied on convolutional architectures.
Representative work such as HRNet \cite{sun2019hrnet} demonstrated that
maintaining high-resolution feature representations throughout the network
improves keypoint localization accuracy, but convolutional
receptive fields remain inherently local, limiting these models' ability
to reason jointly across distant body parts and across long temporal
windows. The introduction of the Transformer architecture
\cite{vaswani2017attention} and its extension to vision
\cite{dosovitskiy2020vit} offered a remedy: self-attention allows every
joint, at every frame, to directly attend to every other joint and frame.
PoseFormer \cite{zheng2021poseformer} was the first to apply spatial and
temporal self-attention to 3D HPE, and subsequent work refined this idea
along several axes. MHFormer \cite{li2022mhformer} models multiple
plausible pose hypotheses to resolve depth ambiguity, MixSTE
\cite{zhang2022mixste} alternates spatial and temporal transformer blocks
in a sequence-to-sequence formulation, PoseFormerV2
\cite{zhao2023poseformerv2} exploits the frequency domain to reduce
sequence length, STCFormer \cite{tang2023stcformer} models spatial and
temporal correlations jointly through criss-cross attention, and
MotionBERT \cite{zhu2023motionbert} and MotionAGFormer
\cite{mehraban2024motionagformer} pursue unified motion representations
and hybrid Transformer-GCN designs, respectively. However, the
self-attention mechanism these methods rely on incurs quadratic
computational and memory complexity in sequence length, and processes
spatio-temporal correlations in a manner that does not naturally
distinguish the directionality or hierarchy inherent to human skeletal
structure. As the field pushes toward longer input sequences for more
temporally consistent predictions, this quadratic cost becomes an
increasingly severe practical bottleneck.

The open question is whether a sequence model can capture the same
long-range spatio-temporal dependencies that make Transformers effective
for 3D HPE while scaling linearly rather than quadratically with sequence
length. Structured state space models (SSMs) have recently emerged as a
candidate answer. The S4 model \cite{gu2021s4} demonstrated that a
structured, continuous-time state space formulation can model long
sequences efficiently, and Mamba \cite{gu2023mamba} extended this line of
work with input-dependent (selective) parameterization and a
hardware-aware scan algorithm, achieving Transformer-competitive quality
with linear-time complexity. This efficiency prompted rapid adoption in
vision: Vision Mamba \cite{zhu2024visionmamba} and VMamba
\cite{liu2024vmamba} adapted the originally causal, unidirectional SSM
formulation to non-causal 2D visual data through bidirectional and
multi-directional scanning strategies, and SSMs have since been applied to
human motion and pose estimation directly
\cite{mondal2024hummuss, zhang2024posemagic, nguyen2025skeletonmamba,
zheng2025mambatopology, lu2025structureawaremamba, feng2025highresssm,
cui2025sasmamba, wang2025dbmambapose}.

Across this emerging body of work, a plain unidirectional SSM applied
naively to skeletal data tends to underperform, and some form of
bidirectional processing or custom scan ordering is needed to compensate
for the fact that an SSM processes its input as a
strict 1D sequence, whereas a human skeleton is not linear but a
hierarchical, tree-structured graph of joints connected by bones
\cite{yan2018stgcn, loper2015smpl}. Existing remedies to this mismatch are
largely heuristic or learned: some methods reorder joints according to an
ad hoc geometric heuristic, while others, such as the Group Scan and Joint
Scan strategies explored in \cite{nguyen2025skeletonmamba}, learn scan
groupings directly from data. Yet the general finding that a principled,
deterministic node ordering can markedly improve sequential modeling of
graph-structured data is well established outside the pose estimation
literature. Breadth-first search (BFS) ordering, for instance, has been
shown to make auto-regressive sequence models over graphs both more
scalable and more effective by exposing local graph structure to the
sequence model in a consistent, canonical way \cite{you2018graphrnn}. This
suggests that the skeleton's own kinematic tree, rather than a learned or
arbitrary heuristic, may offer a more principled scan order for SSM-based
pose estimation.

In this thesis, we investigate this direction directly. We adopt a
bidirectional spatio-temporal state space architecture for monocular 3D
human pose estimation, and introduce a breadth-first search (BFS)
kinematic-tree scan as an architectural refinement to the spatial scan
order: joints are traversed outward from a root joint (the pelvis)
following the parent-child hierarchy defined by the body's kinematic tree
\cite{loper2015smpl}, layer by layer, before being flattened into the 1D
sequence consumed by the state space model. This is motivated by the
hypothesis that a scan order derived directly from skeletal anatomy,
rather than from a fixed joint index or a learned grouping, better
preserves local joint adjacency for the SSM's local modeling capacity
while remaining simple, deterministic, and free of additional learned
parameters. We evaluate this approach on the standard Human3.6M
\cite{ionescu2014human36m} and MPI-INF-3DHP \cite{mehta2017mpiinf3dhp}
benchmarks against established Transformer-, GCN-, and SSM-based
baselines.

\section{Problem Statement}
\label{sec:intro-problem-statement}

The literature reviewed in Section~\ref{sec:intro-overview}
reveals two distinct, compounding problems that motivate the architecture
developed in this thesis.

\textbf{Problem 1: the computational cost of attention-based
spatio-temporal modeling.} Transformer-based 3D HPE methods
\cite{zheng2021poseformer, li2022mhformer, zhang2022mixste,
zhao2023poseformerv2, tang2023stcformer, zhu2023motionbert,
mehraban2024motionagformer} rely on self-attention to model relationships
across joints and frames, but self-attention's computational and memory
cost scales quadratically with input sequence length. Since prediction
accuracy in 3D HPE benefits from longer temporal context (more frames
give the model more evidence to disambiguate depth and resolve occlusion),
this quadratic scaling directly limits how much temporal context these
models can practically exploit, and forces a trade-off between temporal
receptive field and computational budget that becomes increasingly costly
as sequence length grows.

\textbf{Problem 2: a mismatch between sequential state space modeling and
the skeleton's non-sequential structure.} State space models such as
Mamba \cite{gu2023mamba} process input as a strict one-dimensional
sequence, and were originally formulated for inherently sequential or
unidirectional data. A human skeleton, however, is not a sequence but a
hierarchical, tree-structured graph of joints connected by rigid bones
\cite{loper2015smpl}, with no canonical linear order. Flattening this
graph into a 1D sequence therefore requires an explicit choice of scan
order, and this choice determines which joints appear adjacent to one
another in the sequence the SSM processes, which directly affects the
model's ability to capture local joint dependencies. Existing approaches
to this problem are unsatisfying in one of two ways: they either impose a
fixed but anatomically arbitrary order (e.g., dataset joint-index order,
or a hand-designed geometric heuristic), or they learn a scan grouping
from data, such as the Group Scan and Joint Scan strategies in
\cite{nguyen2025skeletonmamba}, at the cost of additional learned
parameters and reduced interpretability. Neither approach derives the scan
order directly and deterministically from the skeleton's own kinematic
hierarchy, despite evidence from outside the pose estimation literature
that principled graph-node orderings, such as breadth-first search, can
markedly improve sequential models over graph-structured data
\cite{you2018graphrnn}.

These two problems together define the central question this
thesis addresses: \textbf{how can a sequence model be designed that
retains the linear-time efficiency of state space models while processing
the human skeleton in a manner that respects its underlying kinematic
structure, rather than an arbitrary or learned joint ordering?} The
remainder of this thesis develops and evaluates an architecture that
answers this question through a bidirectional spatio-temporal state space
backbone combined with a deterministic, breadth-first search
kinematic-tree scan order.

\section{Objectives}
\label{sec:intro-objectives}

Guided by the problem statement in Section~\ref{sec:intro-problem-statement},
the objectives of this thesis are as follows:

\begin{enumerate}
  \item \textbf{To design a bidirectional spatio-temporal state space
    architecture for monocular 3D human pose estimation} that models both
    intra-frame joint relationships and inter-frame temporal correlations
    while retaining the linear-time computational complexity of state
    space models \cite{gu2023mamba, gu2021s4}, avoiding the quadratic
    scaling of self-attention-based approaches \cite{zheng2021poseformer,
    zhang2022mixste, zhu2023motionbert}.

  \item \textbf{To introduce a breadth-first search (BFS) kinematic-tree
    scan order} as the spatial scan strategy for this architecture,
    traversing the skeleton's joints outward from a root joint according
    to the parent-child hierarchy defined by the body's kinematic tree
    \cite{loper2015smpl}, as a deterministic, anatomically-grounded
    alternative to the fixed-index and learned scan orders used in prior
    SSM-based pose estimation work \cite{nguyen2025skeletonmamba},
    motivated by evidence that BFS node orderings improve sequential
    modeling of graph-structured data \cite{you2018graphrnn}.

  \item \textbf{To empirically evaluate the proposed architecture} on the
    standard Human3.6M \cite{ionescu2014human36m} and MPI-INF-3DHP
    \cite{mehta2017mpiinf3dhp} benchmarks, comparing accuracy (MPJPE) and
    computational efficiency (parameter count, MACs) against
    representative Transformer-based \cite{zhang2022mixste,
    zhu2023motionbert, mehraban2024motionagformer}, graph-based
    \cite{yu2023glagcn}, and state-space-based \cite{zhang2024posemagic,
    cui2025sasmamba, wang2025dbmambapose} baselines.

  \item \textbf{To isolate and quantify the specific contribution of the
    BFS kinematic-tree scan} through ablation studies comparing it against
    alternative scan orders (e.g., fixed joint-index order, learned
    group/joint scan strategies), in order to determine whether any
    performance difference is attributable to the scan order itself rather
    than to the underlying SSM backbone alone.
\end{enumerate}

\section{Limitations and Scope}
\label{sec:intro-limitations}

This thesis focuses specifically on the 2D-to-3D lifting stage of the
monocular 3D human pose estimation pipeline. Several boundaries define the
scope of this work:

\begin{itemize}
  \item \textbf{Single-person, monocular RGB input.} The proposed
    architecture is designed and evaluated for single-person 3D pose
    estimation from a single RGB camera viewpoint. Multi-person scenarios,
    multi-view fusion, and depth-sensor or inertial-sensor inputs fall
    outside the scope of this thesis.

  \item \textbf{2D keypoint input assumed.} Consistent with the standard
    2D-to-3D lifting formulation \cite{pavllo2019videopose3d}, this thesis
    assumes 2D keypoints are available as input, either from ground-truth
    annotations or an off-the-shelf 2D pose detector \cite{sun2019hrnet}.
    Improving 2D keypoint detection accuracy itself is not an objective of
    this work; errors introduced at the 2D detection stage are treated as
    an external input condition rather than a problem this architecture
    addresses.

  \item \textbf{Evaluation limited to standard benchmark datasets.} The
    proposed method is evaluated on Human3.6M \cite{ionescu2014human36m}
    and MPI-INF-3DHP \cite{mehta2017mpiinf3dhp}, both of which are
    captured under controlled or semi-controlled conditions. Performance
    on fully unconstrained in-the-wild video, where ground-truth 3D
    annotations are unavailable for direct quantitative evaluation, is not
    assessed.

  \item \textbf{Fixed, known skeletal topology.} The breadth-first search
    kinematic-tree scan relies on a fixed, predefined parent-child joint
    hierarchy \cite{loper2015smpl} consistent with the skeletal
    definitions used by the evaluation datasets. This thesis does not
    address generalization to non-standard skeletal topologies (e.g.,
    pathological anatomy, partial-body input, or rigs with a different
    joint count/hierarchy) that would require re-deriving the traversal
    order.

  \item \textbf{Efficiency measured, not deployed.} Computational
    efficiency is assessed in terms of parameter count and
    multiply-accumulate operations (MACs), consistent with prior work
    \cite{zhang2024posemagic, cui2025sasmamba}. This thesis does not
    include on-device deployment, latency benchmarking on embedded or
    mobile hardware, or real-time system integration.

  \item \textbf{Architectural scope.} The contribution of this thesis is
    confined to the sequence backbone and spatial scan order of the
    lifting network. Improvements to loss functions, data augmentation
    strategies, temporal post-processing/smoothing, or multi-hypothesis
    modeling \cite{li2022mhformer} are outside the scope of this work,
    except where a standard, established choice is adopted without
    modification to support the core architectural comparison.
\end{itemize}

\section{Thesis Outline}
\label{sec:intro-outline}

The remainder of this thesis is organized as follows.

Chapter~\ref{ch:literature-review} reviews the literature relevant to this
work, tracing the evolution of monocular 3D human pose estimation from
convolutional and graph-based methods, through Transformer-based
architectures, to the recent emergence of state space models, and
situates the kinematic-tree scan order proposed in this thesis within that
context.

Chapter~\ref{ch:methodology} presents the proposed architecture in detail,
including the bidirectional spatio-temporal state space backbone and the
breadth-first search kinematic-tree scan.

Chapter~\ref{ch:results} reports the experimental setup and benchmark
results on Human3.6M and MPI-INF-3DHP, comparisons against Transformer-,
graph-, and state-space-based baselines, and ablation studies isolating
the effect of the proposed scan order.

Finally, Chapter~\ref{ch:conclusion} summarizes the contributions of this
thesis, discusses its limitations, and outlines directions for future
work.

\FloatBarrier

